\chapter{Especificación de Requisitos}
\label{cap:requisitos}

\section{Requisitos Funcionales}

\begin{itemize}
    \item \textbf{RF-01. Control de temperatura HVAC:} El sistema debe permitir ajustar el modo de funcionamiento y la temperatura de consigna de la unidad de conductos General/Fujitsu mediante comandos del orquestador. \textbf{[Pendiente de implementar.]}

    \item \textbf{RF-02. Lectura de sensores ambientales:} El sistema debe leer periódicamente la temperatura y la humedad relativa de los sensores Xiaomi BLE y almacenarlos en la base de datos local. \textbf{[Verificado 13/05/2026]} --- El módulo \texttt{ble\_sensors.py} lee el sensor LYWSD03MMC mediante conexión GATT directa, sin flashear firmware ni depender de la nube de Xiaomi.

    \item \textbf{RF-03. Control de dispositivos Tuya:} El sistema debe poder activar y desactivar el enchufe inteligente de forma local, sin requerir conexión a Internet. \textbf{[Verificado 13/05/2026]} --- Respuesta ON/OFF en menos de 2 segundos sobre la red air-gapped mediante \texttt{tinytuya}~\cite{tinytuya}.

    \item \textbf{RF-04. Control del ventilador FanLamp F8808:} El sistema debe poder controlar el ventilador de techo (encendido/apagado, 5 velocidades, luz y modo noche) sin depender de la aplicación del fabricante. \textbf{[Verificado mayo 2026]} --- 11 comandos mapeados mediante anuncios BLE.

    \item \textbf{RF-05. Orquestación por agentes de IA:} El orquestador OpenClaw~\cite{openclaw} debe analizar los datos ambientales y generar recomendaciones o ejecutar acciones de control de forma autónoma. \textbf{[Implementado]} --- Model Router v2 con tres niveles de enrutamiento.

    \item \textbf{RF-06. Sensorización continua:} El sistema debe leer el sensor ambiental cada 60 segundos y mantener un histórico accesible en tiempo real. \textbf{[Implementado]} --- Sensor Daemon como servicio systemd.

    \item \textbf{RF-07. Evaluación de confort térmico:} El sistema debe calcular el índice de calor (NOAA) y asignar zonas de confort. \textbf{[Implementado]} --- Comfort Advisor con 6 niveles térmicos.

    \item \textbf{RF-08. Recomendaciones con contexto exterior:} El sistema debe contrastar las condiciones interiores con las exteriores y recomendar acciones sobre ventanas, persianas y ventilación. \textbf{[Implementado]} --- Smart Advisor.

    \item \textbf{RF-09. Notificaciones proactivas:} El sistema debe notificar al usuario cuando ejecuta acciones automáticas o necesita confirmación. \textbf{[Implementado]} --- Notificaciones vía Telegram cada 2 minutos.

    \item \textbf{RF-10. Auto-recuperación ante fallos:} El sistema debe detectar fallos en los modelos de lenguaje y en la conectividad con los dispositivos, recuperándose automáticamente. \textbf{[Implementado]} --- Circuit breaker, reintentos con backoff, SelfHealManager.

    \item \textbf{RF-11. Registro de analítica energética:} El sistema debe registrar el historial de estados de actuadores y sensores. \textbf{[Parcial]} --- Histórico JSONL del Sensor Daemon. Base de datos SQLite pendiente.

    \item \textbf{RF-12. Interfaz de control por Telegram:} El sistema debe exponer una interfaz por bot de Telegram para consultar el estado y ejecutar comandos. \textbf{[Implementado]} --- Canal nativo del gateway OpenClaw.
\end{itemize}

\section{Requisitos No Funcionales}

\begin{itemize}
    \item \textbf{RNF-01. Funcionamiento 100\% local:} Ninguna función crítica del sistema puede depender de conectividad a Internet. \textbf{[Verificado]} --- El sistema opera sobre un segmento de red air-gapped sin salida a Internet. Todo el procesamiento, desde la inferencia de modelos hasta el control de dispositivos, se ejecuta en el servidor local.

    \item \textbf{RNF-02. Privacidad por diseño:} Los datos de telemetría del hogar no deben transmitirse a servidores de terceros. \textbf{[Verificado]} --- El aislamiento físico de la red impide cualquier fuga de datos. Las \textit{Local Keys} de los dispositivos Tuya se almacenan exclusivamente en el servidor.

    \item \textbf{RNF-03. Baja latencia de control:} Tiempo de respuesta inferior a 2 segundos en condiciones normales de red local. \textbf{[Verificado 13/05/2026]} --- El enchufe Tuya responde en menos de 2~s. Las consultas de estado vía caché LRU se sirven en menos de 1~ms.

    \item \textbf{RNF-04. Resiliencia ante fallos:} Un fallo en el módulo de IA no debe dejar la unidad HVAC en estado de error. \textbf{[Parcial]} --- El \textit{circuit breaker} del Model Router y el SelfHealManager están implementados y verificados. La parte específica del firmware del ESP32 (mantener último estado válido en el bus LIN) está pendiente de desarrollar.

    \item \textbf{RNF-05. Extensibilidad:} La arquitectura debe permitir agregar nuevos tipos de dispositivos sin modificar el núcleo del orquestador. \textbf{[Verificado]} --- El sistema integra tres protocolos distintos (Tuya por Wi-Fi local, Xiaomi por BLE GATT, FanLamp por BLE advertisements) mediante \textit{drivers} Python independientes, sin modificar el orquestador OpenClaw ni el Model Router.
\end{itemize}
